Este Trabalho de Conclusão de Curso II consolida a avaliação de uma plataforma robótica física para execução de SLAM visual embarcado, deslocando a ênfase de simulações para a análise das restrições reais de hardware. O protótipo desenvolvido, baseado em Raspberry Pi 4, ROS 2 Jazzy e câmera RGB-D Kinect v1, demonstrou ser um artefato funcional para integrar percepção, controle e preparação do mapeamento sob restrições de baixo custo.

As evidências levantadas até o momento indicam que a viabilidade do SLAM em hardware limitado não é condicionada apenas pela complexidade algorítmica, mas principalmente pela estabilidade da cadeia de aquisição de dados. O principal gargalo identificado foi a instabilidade na aquisição sincronizada de fluxos RGB-D, que atua como um gate técnico crítico para a validação do mapeamento. Observou-se que a carga no barramento USB e na CPU durante a aquisição de imagens pode comprometer a sincronia dos dados, exigindo novos ensaios e ajustes antes de concluir a avaliação do RTAB-Map.

Apesar dessas limitações, o trabalho atingiu seus objetivos de estruturar um pipeline de percepção, controle e localização a ser validado progressivamente em ambiente interno real. A integração entre o microcontrolador Arduino Nano e o Raspberry Pi via ROS 2 forneceu uma base de odometria, telemetria e transformações (TF), fundamentais para os próximos ensaios de pose estimada e construção de mapa.

Como contribuição, este estudo fornece um roteiro técnico e metodológico para o desenvolvimento de robôs de serviço acessíveis, mapeando os desafios práticos que surgem ao levar o SLAM visual do ambiente simulado para o hardware embarcado. Para trabalhos futuros, sugere-se a integração da pilha de navegação completa (\textit{Nav2}) e a investigação de métodos de compressão, redução de resolução ou pré-processamento de imagens em hardware para mitigar a instabilidade sensorial observada.
